\chapter{Mô tả tổng quan hệ thống}

\section{Tổng quan hệ thống}

Hệ thống Quản lý dịch vụ bảo trì máy lọc nước là một ứng dụng Web được xây dựng nhằm hỗ trợ doanh nghiệp quản lý toàn bộ thông tin liên quan đến sản phẩm, thiết bị, khách hàng, hợp đồng bảo trì, tồn kho, nhân viên và các hoạt động phát sinh trong quá trình vận hành.

Hệ thống được thiết kế theo mô hình quản lý tập trung, cho phép người dùng truy cập thông qua trình duyệt Web để thực hiện các nghiệp vụ tương ứng với quyền hạn được cấp. Bên cạnh các chức năng quản lý thông tin, hệ thống còn hỗ trợ hiển thị Dashboard thống kê và theo dõi dữ liệu hoạt động của từng thiết bị thông qua các biểu đồ trực quan.

Phiên bản hiện tại được phát triển với mục đích xây dựng một hệ thống Demo phục vụ quá trình thực tập, đồng thời làm nền tảng cho việc mở rộng sang hệ thống quản lý thực tế trong tương lai.

\subsection*{Backend Integration}


Các module chính của hệ thống bao gồm:

\begin{itemize}
    \item Dashboard
    \item Quản lý sản phẩm (Product)
    \item Quản lý tồn kho (Inventory)
    \item Quản lý lô hàng (Batch)
    \item Quản lý khách hàng (Customer)
    \item Quản lý hợp đồng (Contract)
    \item Quản lý thiết bị (Device)
    \item Dashboard theo dõi thiết bị (Device Dashboard)
    \item Quản lý nhân viên (Employee)
    \item Activity Log
    \item Hồ sơ cá nhân (Profile)
    \item Đăng nhập và phân quyền người dùng
\end{itemize}

Hình \ref{fig:system_overview} minh họa kiến trúc tổng quan của hệ thống.

\begin{figure}[H]
    \centering
    \includegraphics[width=0.9\textwidth]{figures/architecture/system_overview.png}
    \caption{Kiến trúc tổng quan của hệ thống}
    \label{fig:system_overview}
\end{figure}

%%%%%%%%%%%%%%%%%%%%%%%%%%%%%%%%%%%%%%%%%%%%%%%%%%

\section{Góc nhìn sản phẩm (Product Perspective)}

Hệ thống được phát triển như một ứng dụng Web tích hợp (Integrated Monolith) theo mô hình Client -- Server, nơi Frontend (giao diện Blade) và Backend (REST API) được triển khai trong cùng một ứng dụng Laravel.

\subsection*{Kiến trúc Integrated Monolith}

\begin{itemize}
    \item \textbf{Frontend}: Blade Templates + Bootstrap 5 + JavaScript
    \item \textbf{Backend}: Laravel Controllers + Services + Eloquent ORM
    \item \textbf{API Layer}: REST API với 30+ endpoints JSON
    \item \textbf{Database}: MySQL với 17 tables và relationships
\end{itemize}

Người dùng truy cập hệ thống thông qua trình duyệt Web. Mọi yêu cầu từ giao diện người dùng sẽ được gửi đến Web Server, nơi ứng dụng Laravel xử lý nghiệp vụ thông qua Controllers và Services, thực hiện các thao tác với cơ sở dữ liệu MySQL thông qua Eloquent ORM, sau đó trả kết quả về giao diện. Bên cạnh đó, hệ thống còn cung cấp REST API cho phép các ứng dụng bên ngoài truy vấn và cập nhật dữ liệu.

Để xây dựng giao diện quản trị hiện đại và đồng nhất, hệ thống sử dụng Bootstrap 5 kết hợp với Blade Template Engine nhằm tăng khả năng Responsive và cải thiện trải nghiệm người dùng.

\subsection*{Backend Stack}

\begin{itemize}
    \item \textbf{Framework}: Laravel 13
    \item \textbf{ORM}: Eloquent (17 Models)
    \item \textbf{Database}: MySQL 5.7+ / MariaDB
    \item \textbf{API}: REST endpoints (JSON responses)
    \item \textbf{Validation}: Form Requests
    \item \textbf{Business Logic}: Service Layer
\end{itemize}

Trong định hướng phát triển tương lai, hệ thống có thể tích hợp với nền tảng IoT để tiếp nhận dữ liệu từ các thiết bị máy lọc nước thông qua giao thức MQTT. Các dữ liệu này sẽ được lưu trữ trong cơ sở dữ liệu và hiển thị dưới dạng Dashboard theo từng thiết bị. Ngoài ra, hệ thống cũng hỗ trợ mở rộng thêm các tính năng như Authentication (Sanctum), Authorization (RBAC), và caching (Redis).

Kiến trúc tổng thể của hệ thống được minh họa trong Hình \ref{fig:product_perspective}.

\begin{figure}[H]
    \centering
    \includegraphics[width=\textwidth]{figures/architecture/software_architecture.png}
    \caption{Kiến trúc tổng thể của hệ thống}
    \label{fig:product_perspective}
\end{figure}

%%%%%%%%%%%%%%%%%%%%%%%%%%%%%%%%%%%%%%%%%%%%%%%%%%

\section{Các nhóm người dùng (User Roles)}

Hệ thống được thiết kế với cơ chế phân quyền theo vai trò nhằm đảm bảo mỗi người dùng chỉ được phép sử dụng các chức năng phù hợp với nhiệm vụ của mình.

\subsection{Quản trị viên (Administrator)}

Administrator là người có toàn quyền quản trị hệ thống.

Các chức năng chính bao gồm:

\begin{itemize}
    \item Quản lý tài khoản người dùng.
    \item Quản lý sản phẩm.
    \item Quản lý tồn kho.
    \item Quản lý lô hàng.
    \item Quản lý khách hàng.
    \item Quản lý hợp đồng.
    \item Quản lý thiết bị.
    \item Quản lý nhân viên.
    \item Theo dõi Activity Log.
    \item Xem Dashboard thống kê.
\end{itemize}

\subsection{Nhân viên (Employee)}

Nhân viên chịu trách nhiệm thực hiện các nghiệp vụ hằng ngày.

Bao gồm:

\begin{itemize}
    \item Quản lý khách hàng.
    \item Quản lý hợp đồng.
    \item Quản lý thiết bị.
    \item Quản lý tồn kho.
    \item Theo dõi Dashboard.
\end{itemize}

\subsection{Kỹ thuật viên (Technician)}

Kỹ thuật viên phụ trách việc lắp đặt và bảo trì thiết bị.

Các chức năng bao gồm:

\begin{itemize}
    \item Xem danh sách thiết bị được phân công.
    \item Cập nhật thông tin bảo trì.
    \item Theo dõi Dashboard của thiết bị.
    \item Xem lịch sử bảo trì.
\end{itemize}

Bảng \ref{tab:user_roles} mô tả quyền truy cập của từng nhóm người dùng.

\begin{table}[H]
\centering
\caption{Phân quyền người dùng}
\label{tab:user_roles}

\begin{tabularx}{\textwidth}{|l|c|c|c|}
\hline
\textbf{Chức năng} & \textbf{Admin} & \textbf{Employee} & \textbf{Technician} \\
\hline
Dashboard & \checkmark & \checkmark & \checkmark \\
\hline
Quản lý sản phẩm & \checkmark & & \\
\hline
Quản lý tồn kho & \checkmark & \checkmark & \\
\hline
Quản lý lô hàng & \checkmark & & \\
\hline
Quản lý khách hàng & \checkmark & \checkmark & \\
\hline
Quản lý hợp đồng & \checkmark & \checkmark & \\
\hline
Quản lý thiết bị & \checkmark & \checkmark & \checkmark \\
\hline
Dashboard thiết bị & \checkmark & \checkmark & \checkmark \\
\hline
Quản lý nhân viên & \checkmark & & \\
\hline
Activity Log & \checkmark & & \\
\hline
\end{tabularx}

\end{table}

%%%%%%%%%%%%%%%%%%%%%%%%%%%%%%%%%%%%%%%%%%%%%%%%%%

\section{Môi trường hoạt động (Operating Environment)}

Hệ thống được phát triển theo kiến trúc Web và có thể triển khai trên các máy chủ hỗ trợ PHP.

\subsection{Môi trường phía người dùng (Client)}

Người dùng truy cập hệ thống thông qua các trình duyệt phổ biến như:

\begin{itemize}
    \item Google Chrome
    \item Microsoft Edge
    \item Mozilla Firefox
\end{itemize}

Giao diện được thiết kế theo Responsive Design nhằm hỗ trợ nhiều kích thước màn hình khác nhau.

\subsection{Môi trường phía máy chủ (Server)}

Máy chủ triển khai hệ thống yêu cầu các thành phần sau:

\begin{itemize}
    \item Hệ điều hành Windows hoặc Linux.
    \item Apache Web Server.
    \item PHP 8.x.
    \item MySQL 8.x.
\end{itemize}

\subsection{Môi trường phát triển}

Bảng \ref{tab:development_environment} mô tả các công nghệ sử dụng trong quá trình phát triển hệ thống.

\begin{table}[H]
\centering
\caption{Môi trường phát triển}
\label{tab:development_environment}

\begin{tabularx}{\textwidth}{|l|X|}
\hline
\textbf{Thành phần} & \textbf{Công nghệ sử dụng} \\
\hline
Backend Framework & Laravel 13 \\
\hline
Ngôn ngữ lập trình & PHP 8.5+, HTML5, CSS3, JavaScript \\
\hline
Template Engine & Blade (Laravel) \\
\hline
Framework CSS & Bootstrap 5 \\
\hline
ORM & Eloquent (Laravel) \\
\hline
Hệ quản trị cơ sở dữ liệu & MySQL 5.7+ / MariaDB \\
\hline
Web Server & Apache / Laravel Artisan Server \\
\hline
Development Server & XAMPP / Laravel Artisan \\
\hline
IDE & Visual Studio Code \\
\hline
API Testing & cURL / Postman \\
\hline
Quản lý dependencies & Composer \\
\hline
Quản lý mã nguồn & Git, GitHub \\
\hline
\end{tabularx}

\end{table}

